\section{Estrategias de experimentación}

En las unidades anteriores exploramos los métodos inferenciales para estimar parámetros y contrastar hipótesis relativas a una o dos poblaciones. Sin embargo, en la investigación científica contemporánea, la ingeniería de sistemas, la biotecnología y la ciencia de datos, rara vez nos limitamos a comparar únicamente dos escenarios. Frecuentemente necesitamos evaluar múltiples algoritmos de aprendizaje automático de manera simultánea, comparar varias arquitecturas de servidores web, optimizar campañas de marketing digital con diversas configuraciones multivariadas (pruebas A/B/n) o evaluar el rendimiento de distintos tratamientos bajo condiciones experimentales controladas.

Intentar resolver un problema de comparación entre $k$ grupos realizando pruebas de hipótesis $t$ de Student dos a dos ($\binom{k}{2}$ combinaciones) es un error metodológico crítico, ya que produce una \textbf{inflación severa de la tasa de error Tipo I global} (probabilidad de cometer al menos un falso positivo en la familia de comparaciones). Por ejemplo, si comparamos $k = 5$ algoritmos ($10$ pares posibles) utilizando un nivel de significancia $\alpha = 0.05$ para cada prueba independiente, la probabilidad de encontrar al menos una diferencia "significativa" por puro azar cuando en realidad todos los algoritmos son idénticos asciende a:
\begin{align}
	\alpha_{\text{global}} = 1 - (1 - \alpha)^m = 1 - (1 - 0.05)^{10} = 1 - (0.95)^{10} \approx 0.4013 \quad (40.13\%).
\end{align}
Para evitar este sesgo sistemático y extraer conclusiones válidas sobre múltiples grupos en un solo paso analítico, recurrimos a la teoría de \textbf{Diseño de Experimentos (DoE)} y a su herramienta metodológica primordial: el \textbf{Análisis de Varianza (ANOVA)}, desarrollado originalmente por Sir Ronald A. Fisher.

\subsection{Principios fundamentales y terminología del diseño experimental}

Un diseño experimental riguroso es la planificación sistemática de un estudio de manera que los datos recolectados puedan ser analizados mediante métodos estadísticos objetivos, permitiendo deducir relaciones de causa y efecto con la máxima precisión y el mínimo costo de muestreo.

\begin{definicion}[Terminología formal en el Diseño de Experimentos]
	\begin{itemize}
		\item \textbf{Unidad experimental:} Es la entidad física o informacional básica sobre la cual se aplica un tratamiento o se realiza una medición (por ejemplo, un servidor en la nube, un paciente, un lote de producción o una consulta de base de datos).
		\item \textbf{Factor (Variable independiente):} Es una variable controlada explícitamente por el investigador para evaluar su efecto sobre el proceso (por ejemplo, el tipo de algoritmo, el nivel de optimización del compilador o la dosis de un fármaco).
		\item \textbf{Nivel del factor:} Cada uno de los valores, categorías o modalidades específicas que puede adoptar un factor en el experimento.
		\item \textbf{Tratamiento:} Es una combinación específica de niveles de uno o varios factores aplicada a las unidades experimentales. En un experimento de un solo factor, los tratamientos son sinónimos de los niveles del factor.
		\item \textbf{Variable respuesta (Variable dependiente):} Es la característica o métrica numérica observada y registrada tras la aplicación del tratamiento (por ejemplo, tiempo de latencia en milisegundos, exactitud o \emph{accuracy} de un clasificador, o rendimiento en transacciones por segundo).
		\item \textbf{Error experimental (Ruido):} Representa la variabilidad aleatoria intrínseca y no controlada entre observaciones correspondientes a unidades experimentales tratadas de manera idéntica.
	\end{itemize}
\end{definicion}

La validez y eficiencia del Diseño de Experimentos descansa sobre \textbf{tres principios cardinales enunciados por Fisher}:

\begin{teorema}[Los Tres Principios de Fisher en el DoE]
	\begin{enumerate}
		\item \textbf{Replicación (\emph{Replication}):} Es la repetición de un mismo tratamiento sobre múltiples unidades experimentales independientes. Cumple dos funciones vitales: (a) proporciona una estimación interna y objetiva de la varianza del error experimental ($\sigma^2$), sin la cual es imposible evaluar la significancia estadística de los tratamientos; y (b) reduce el error estándar de la media del tratamiento ($\sigma / \sqrt{n}$), aumentando la precisión de la estimación.
		\item \textbf{Aleatorización (\emph{Randomization}):} Es la asignación estrictamente probabilística (al azar) de las unidades experimentales a los tratamientos y del orden temporal/espacial de las pruebas. Garantiza la independencia de los errores ($\epsilon_{ij}$), elimina sesgos sistemáticos del operador u orden de ejecución, y valida las distribuciones de probabilidad utilizadas por las pruebas inferenciales ($F$ y $t$).
		\item \textbf{Control Local o Bloqueo (\emph{Blocking}):} Es una técnica de agrupación que consiste en dividir el conjunto de unidades experimentales en subgrupos o "bloques" lo más homogéneos posible entre sí, aislando en una dimensión separada la varianza producida por fuentes de ruido conocidas (como diferentes lotes de hardware, turnos de trabajo o configuraciones del sistema operacionales). El bloqueo incrementa de forma contundente la potencia de la prueba al reducir la varianza residual del error experimental.
	\end{enumerate}
\end{teorema}

\section{Experimentos con un factor: análisis de varianza}

El \textbf{Análisis de Varianza de un factor} (o ANOVA completamente aleatorizado) se utiliza cuando deseamos comparar las medias poblacionales de $k$ tratamientos (o niveles de un único factor) habiendo asignado al azar $N$ unidades experimentales de forma que cada tratamiento $i$ cuenta con $n_i$ réplicas ($N = \sum_{i=1}^k n_i$).

\subsection{Modelo estadístico lineal paramétrico}

Toda observación experimental $Y_{ij}$ (la $j$-ésima réplica bajo el $i$-ésimo tratamiento) se puede expresar formalmente mediante el siguiente modelo paramétrico aditivo:
\begin{align}
	Y_{ij} = \mu + \tau_i + \epsilon_{ij}, \quad \text{para } i = 1, 2, \ldots, k \text{ y } j = 1, 2, \ldots, n_i,
	\label{eq:modelo_anova1}
\end{align}
donde:
\begin{itemize}
	\item $\mu$ es la media general o global de toda la población de observaciones.
	\item $\tau_i = \mu_i - \mu$ es el \textbf{efecto específico del tratamiento $i$}, que mide cuánto se desvía la media real del tratamiento $i$ ($\mu_i$) respecto al promedio general $\mu$. Por construcción matemática, la suma ponderada de los efectos de los tratamientos es nula: $\sum_{i=1}^k n_i \tau_i = 0$.
	\item $\epsilon_{ij}$ es el término del error experimental aleatorio asociado a la observación $(i,j)$, el cual se asume distribuido de forma normal, independiente e idénticamente distribuida con media cero y varianza constante $\sigma^2$:
	\begin{align}
		\epsilon_{ij} \sim N(0, \sigma^2) \quad \text{i.i.d.}
	\end{align}
\end{itemize}

Bajo esta formulación, el objetivo fundamental del ANOVA de un factor es contrastar la hipótesis nula de que ningún tratamiento tiene un efecto diferencial sobre la variable respuesta:
\begin{align}
	H_0&: \tau_1 = \tau_2 = \ldots = \tau_k = 0 \quad \Longleftrightarrow \quad \mu_1 = \mu_2 = \ldots = \mu_k, \\
	H_a&: \tau_i \neq 0 \text{ para al menos un tratamiento } i \quad (\text{al menos una media poblacional difiere}).
\end{align}

\subsection{Descomposición fundamental de la suma de cuadrados}

La brillante conceptualización de Fisher para resolver el problema comparativo consiste en \textbf{particionar la varianza total} del conjunto de datos en dos componentes aditivos perfectamente diferenciados: la variabilidad atribuible a las diferencias \emph{entre} los tratamientos y la variabilidad debida al error experimental aleatorio \emph{dentro} de cada tratamiento.

\begin{teorema}[Teorema de partición de la suma de cuadrados en ANOVA de un factor]
	Sea $\bar{Y}_{i.} = \frac{1}{n_i}\sum_{j=1}^{n_i} Y_{ij}$ la media muestral del tratamiento $i$, y sea $\bar{Y}_{..} = \frac{1}{N}\sum_{i=1}^k \sum_{j=1}^{n_i} Y_{ij}$ la media general de todas las $N$ observaciones. La \textbf{Suma de Cuadrados Total ($\text{SCT}$)} del experimento satisface la siguiente igualdad exacta:
	\begin{align}
		\sum_{i=1}^k \sum_{j=1}^{n_i} (Y_{ij} - \bar{Y}_{..})^2 = \sum_{i=1}^k n_i (\bar{Y}_{i.} - \bar{Y}_{..})^2 + \sum_{i=1}^k \sum_{j=1}^{n_i} (Y_{ij} - \bar{Y}_{i.})^2,
		\label{eq:particion_suma_cuadrados}
	\end{align}
	que expresamos abreviadamente como:
	\begin{align}
		\text{SCT} = \text{SCTR} + \text{SCE},
	\end{align}
	donde:
	\begin{itemize}
		\item $\text{SCT} = \sum_{i=1}^k \sum_{j=1}^{n_i} (Y_{ij} - \bar{Y}_{..})^2$ es la Suma de Cuadrados Total (mide la dispersión global de todos los datos respecto al gran promedio, con $\nu_T = N - 1$ grados de libertad).
		\item $\text{SCTR} = \sum_{i=1}^k n_i (\bar{Y}_{i.} - \bar{Y}_{..})^2$ es la \textbf{Suma de Cuadrados del Tratamiento} (mide la dispersión entre las medias de los grupos y la media general ponderada por sus tamaños muestrales, con $\nu_{\text{TR}} = k - 1$ grados de libertad).
		\item $\text{SCE} = \sum_{i=1}^k \sum_{j=1}^{n_i} (Y_{ij} - \bar{Y}_{i.})^2 = \sum_{i=1}^k (n_i - 1) S_i^2$ es la \textbf{Suma de Cuadrados del Error o residual} (mide la dispersión interna de las observaciones respecto a la media de su propio tratamiento, con $\nu_E = N - k$ grados de libertad).
	\end{itemize}
\end{teorema}

\section{Análisis de efectos de modelo fijo}

Al dividir cada suma de cuadrados entre sus respectivos grados de libertad se obtienen los \textbf{Cuadrados Medios ($\text{CM}$)}, que representan varianzas muestrales insesgadas bajo diferentes condiciones poblacionales:
\begin{align}
	\text{CMTR} &= \frac{\text{SCTR}}{k - 1}, \\
	\text{CME} &= \frac{\text{SCE}}{N - k}.
\end{align}

Para comprender la lógica del estadístico de prueba del ANOVA, debemos analizar la \textbf{esperanza matemática de los cuadrados medios}:
\begin{align}
	E(\text{CME}) &= \sigma^2, \\
	E(\text{CMTR}) &= \sigma^2 + \frac{\sum_{i=1}^k n_i \tau_i^2}{k - 1}.
	\label{eq:esperanza_cmtr}
\end{align}
La ecuación (\ref{eq:esperanza_cmtr}) es reveladora:
\begin{itemize}
	\item Si la hipótesis nula $H_0$ es \textbf{verdadera} ($\tau_i = 0$ para todo $i$), entonces el segundo término se anula y tenemos que $E(\text{CMTR}) = \sigma^2$. En consecuencia, tanto $\text{CMTR}$ como $\text{CME}$ son estimadores independientes de la misma varianza poblacional $\sigma^2$, y su cociente rondará teóricamente el valor $1.0$.
	\item Si la hipótesis nula es \textbf{falsa} (existen tratamientos con $\tau_i \neq 0$), el término $\frac{\sum n_i \tau_i^2}{k-1}$ será estrictamente positivo ($\sum n_i \tau_i^2 > 0$). Por tanto, $E(\text{CMTR}) > \sigma^2$, provocando que el cociente $\text{CMTR} / \text{CME}$ sea significativamente mayor a $1.0$.
\end{itemize}

\begin{teorema}[Estadístico de prueba $F$ en ANOVA]
	Bajo los supuestos del modelo paramétrico y suponiendo verdadera la hipótesis nula $H_0: \tau_i = 0$, el estadístico de prueba:
	\begin{align}
		F = \frac{\text{CMTR}}{\text{CME}} = \frac{\text{SCTR} / (k - 1)}{\text{SCE} / (N - k)}
	\end{align}
	sigue una distribución $F$ de Fisher-Snedecor con $\nu_1 = k - 1$ grados de libertad en el numerador y $\nu_2 = N - k$ grados de libertad en el denominador.
	
	Rechazamos la hipótesis nula $H_0$ al nivel de significancia $\alpha$ si y solo si el valor observado del estadístico excede el cuantil crítico superior de la distribución:
	\begin{align}
		F > F_{\alpha, \, k-1, \, N-k}.
	\end{align}
\end{teorema}

Todos los cálculos y componentes computacionales del procedimiento se sintetizan de forma clara y estandarizada en la \textbf{Tabla de ANOVA de un factor} (Tabla \ref{tab:anova_un_factor}).

\begin{table}[h!]
	\centering
	\begin{tabular}{lcccc}
		\hline
		\textbf{Fuente de Variación} & \textbf{Suma de Cuadrados (SC)} & \textbf{g.l. ($\nu$)} & \textbf{Cuadrado Medio (CM)} & \textbf{Estadístico $F$} \\ \hline
		Tratamientos (Entre) & $\text{SCTR} = \sum n_i (\bar{Y}_{i.} - \bar{Y}_{..})^2$ & $k - 1$ & $\text{CMTR} = \frac{\text{SCTR}}{k - 1}$ & $F = \frac{\text{CMTR}}{\text{CME}}$ \\ [1.5ex]
		Error (Dentro) & $\text{SCE} = \sum (n_i - 1) S_i^2$ & $N - k$ & $\text{CME} = \frac{\text{SCE}}{N - k}$ & \\ [1.5ex] \hline
		\textbf{Total} & $\text{SCT} = \sum \sum (Y_{ij} - \bar{Y}_{..})^2$ & $N - 1$ & & \\ \hline
	\end{tabular}
	\caption{Tabla general del Análisis de Varianza de un factor (ANOVA completamente aleatorizado).}
	\label{tab:anova_un_factor}
\end{table}

\section*{Comparaciones múltiples y pruebas post-hoc}

Cuando la prueba global $F$ del ANOVA es estadísticamente significativa y rechazamos $H_0$, sabemos con certeza inferencial que \textbf{al menos un tratamiento difiere en media de los demás}, pero la tabla ANOVA por sí sola no especifica \emph{cuáles} pares de tratamientos son los responsables de dicha significancia. Para identificar los pares exactos de diferencias sin incurrir en la explosión de falsos positivos ($\alpha$ global), debemos aplicar procedimientos post-hoc de comparación múltiple.

\subsection*{Diferencia Mínima Significativa de Fisher (LSD)}

El procedimiento \textbf{LSD de Fisher (\emph{Least Significant Difference})} es la extensión natural de la prueba $t$ de Student para comparar dos medias grupales post-ANOVA. En lugar de utilizar la varianza combinada particular de los dos grupos bajo comparación ($S_p^2$), el procedimiento aprovecha toda la información del experimento utilizando la estimación global de la varianza del error obtenida del ANOVA ($\text{CME}$ con $N-k$ grados de libertad).

\begin{teorema}[Criterio LSD de Fisher]
	Dos medias muestrales de tratamientos ($\bar{Y}_{i.}$ y $\bar{Y}_{j.}$) con tamaños muestrales $n_i$ y $n_j$ se declaran significativamente diferentes al nivel $\alpha$ si la diferencia absoluta entre ellas excede la Diferencia Mínima Significativa:
	\begin{align}
		|\bar{Y}_{i.} - \bar{Y}_{j.}| > \text{LSD} = t_{\alpha/2, \, N-k} \sqrt{\text{CME} \left( \frac{1}{n_i} + \frac{1}{n_j} \right)},
	\end{align}
	donde $t_{\alpha/2, N-k}$ es el cuantil superior de la distribución $t$ de Student con los grados de libertad del error del ANOVA ($N-k$).
\end{teorema}

\begin{observacion}
	El método LSD es potente (baja probabilidad de error Tipo II), pero \textbf{no controla la tasa de error experimental por familia (FWER)} ante un número elevado de tratamientos ($k > 3$). Por ello, los científicos de datos y experimentadores rigurosos prefieren el método HSD de Tukey.
\end{observacion}

\subsection*{Diferencia Honestamente Significativa de Tukey (HSD)}

El procedimiento \textbf{HSD de Tukey (\emph{Honestly Significant Difference})} fue diseñado específicamente para comparar simultáneamente todos los pares posibles de tratamientos garantizando que la probabilidad global de cometer al menos un falso positivo en todo el conjunto de comparaciones sea estrictamente menor o igual a $\alpha$ ($\text{FWER} \leq \alpha$).

El método utiliza la distribución del \textbf{rango estandarizado ($q$)}, que modela la distribución muestral de la diferencia entre el valor máximo y mínimo de un conjunto de $k$ medias muestrales independientes e idénticamente distribuidas.

\begin{teorema}[Criterio HSD de Tukey para tamaños de muestra iguales]
	Si cada uno de los $k$ tratamientos cuenta con exactamente el mismo número de réplicas ($n_1 = n_2 = \ldots = n_k = n$), dos tratamientos $i$ y $j$ difieren significativamente al nivel de significancia global $\alpha$ si:
	\begin{align}
		|\bar{Y}_{i.} - \bar{Y}_{j.}| > \text{HSD} = q_{\alpha, \, k, \, N-k} \sqrt{\frac{\text{CME}}{n}},
		\label{eq:tukey_hsd}
	\end{align}
	donde $q_{\alpha, k, N-k}$ es el valor crítico en las tablas del rango estandarizado para $k$ tratamientos y $N-k$ grados de libertad en el error.
\end{teorema}

\begin{teorema}[Extensión de Tukey-Kramer para muestras de tamaño desigual]
	Cuando los tamaños muestrales no son balanceados ($n_i \neq n_j$), la generalización rigurosa de Kramer adapta el criterio reemplazando el tamaño muestral constante por la media armónica de los tamaños del par en comparación:
	\begin{align}
		|\bar{Y}_{i.} - \bar{Y}_{j.}| > \text{HSD}_{ij} = \frac{q_{\alpha, \, k, \, N-k}}{\sqrt{2}} \sqrt{\text{CME} \left( \frac{1}{n_i} + \frac{1}{n_j} \right)}.
	\end{align}
\end{teorema}

\subsection*{Corrección de Bonferroni}

La \textbf{corrección de Bonferroni} es un método post-hoc universal, simple y conservador basado en la desigualdad de Boole para probabilidades de eventos intersecados. Si el investigador planea realizar $m = \binom{k}{2}$ comparaciones por pares y desea mantener el nivel de confianza familiar global en al menos $(1-\alpha)100\%$, el nivel de significancia individual para cada prueba particular par a par se ajusta dividiendo el nivel nominal entre el número total de comparaciones:
\begin{align}
	\alpha^* = \frac{\alpha}{m} = \frac{\alpha}{\binom{k}{2}}.
\end{align}
Dos medias se declaran significativamente diferentes si $|t_{ij}| > t_{\alpha^*/2, N-k}$, o bien si su diferencia supera el margen crítico $t_{\alpha/(2m), N-k} \sqrt{\text{CME}(\frac{1}{n_i}+\frac{1}{n_j})}$.

\section{Adecuación del modelo: análisis de residuos}

La validez teórica y la robustez inferencial de cualquier conclusión extraída de un ANOVA descansan fundamentalmente en que los datos experimentales y los residuos ($\hat{\epsilon}_{ij} = Y_{ij} - \hat{Y}_{ij}$) cumplan con los supuestos subyacentes del modelo paramétrico. Todo científico de datos debe auditar y verificar rigurosamente estos tres pilares antes de presentar resultados experimentales en un entorno productivo o de investigación:

\begin{enumerate}
	\item \textbf{Homocedasticidad (Igualdad de varianzas intergrupos):} La varianza del error experimental en todos los $k$ tratamientos debe ser estadísticamente equivalente ($\sigma_1^2 = \sigma_2^2 = \ldots = \sigma_k^2 = \sigma^2$). El ANOVA de un factor es razonablemente robusto a desviaciones moderadas de homoscedasticidad \emph{únicamente si los tamaños muestrales son exactamente iguales ($n_1 = n_2 = \ldots = n_k$)}. Sin embargo, si los tamaños son desiguales, el estadístico $F$ puede volverse altamente impreciso.

	Para contrastar formalmente la hipótesis de homoscedasticidad se emplean dos pruebas de referencia:
	\begin{itemize}
		\item \textbf{Prueba de Bartlett:} Es altamente sensible a la asunción de normalidad de la población. Contraste $H_0: \sigma_1^2 = \ldots = \sigma_k^2$ utilizando un estadístico basado en el cociente ponderado de las varianzas muestrales que sigue una distribución chi-cuadrada con $k-1$ grados de libertad.
		\item \textbf{Prueba de Levene y Brown-Forsythe:} Es el estándar de oro moderno en analítica computacional (integrado en \texttt{scipy.stats.levene}) debido a su extraordinaria robustez ante distribuciones no normales o sesgadas. Consiste en realizar un ANOVA univariado aplicándolo directamente a las desviaciones absolutas de las observaciones respecto a la media (Levene: $Z_{ij} = |Y_{ij} - \bar{Y}_{i.}|$) o respecto a la mediana del tratamiento (Brown-Forsythe: $Z_{ij} = |Y_{ij} - \tilde{Y}_{i.}|$).
	\end{itemize}
	Si la hipótesis de homoscedasticidad es rechazada al nivel $\alpha$, debemos abandonar el ANOVA paramétrico estándar y utilizar el \textbf{ANOVA de Welch sobre una vía} o transformaciones estabilizadoras de varianza de la familia Box-Cox (como $\log Y$ o $\sqrt{Y}$).

	\item \textbf{Normalidad de los errores residuales ($\epsilon_{ij} \sim N(0,\sigma^2)$):} La distribución de los residuos dentro de cada tratamiento o del vector general de residuos del modelo debe aproximarse a una distribución normal. Este supuesto se audita cualitativamente mediante un gráfico de probabilidad normal \textbf{Q-Q Plot} de los residuos, y cuantitativamente mediante la \textbf{Prueba de Shapiro-Wilk} o la \textbf{Prueba de Kolmogorov-Smirnov} sobre los residuos estandarizados.

	Si se detectan violaciones severas a la normalidad (especialmente asimetría extrema u observaciones atípicas extremas con muestras pequeñas), el científico de datos debe utilizar la alternativa no paramétrica del ANOVA: la \textbf{Prueba de Kruskal-Wallis} (basada en rangos y medianas en lugar de medias y varianzas).

	\item \textbf{Independencia de las observaciones:} Los residuos de cualquier par de unidades experimentales deben ser estadísticamente independientes ($Cov(\epsilon_{ij}, \epsilon_{lm}) = 0$). Este es el supuesto más crítico de todos: \textbf{la violación a la independencia infesta directamente la estimación del Cuadrado Medio del Error ($\text{CME}$)}, invalidando por completo la prueba $F$. La independencia no se puede arreglar fácilmente a posteriori con transformaciones algebraicas; se garantiza a priori mediante la estricta \textbf{aleatorización en el diseño de muestreo} y la correcta identificación de dependencias temporales u espaciales (que requerirían modelos de series de tiempo o modelos de efectos mixtos longitudinales).
\end{enumerate}

\section{Bloques completamente aleatorizados, cuadrados latinos y grecolatinos}

En innumerables escenarios prácticos, las unidades experimentales disponibles no son un conjunto estrictamente homogéneo. Por ejemplo, al evaluar motores de bases de datos en servidores físicos de diversas generaciones, o al probar fármacos en pacientes de distintos rangos de edad, una gran porción de la variabilidad observada en la variable respuesta no se debe al factor bajo estudio (el tratamiento), sino a las diferencias intrínsecas de las unidades experimentales.

Si realizáramos un ANOVA de un factor ignorando esa heterogeneidad conocida, dicha variación se absorbería directamente dentro de la Suma de Cuadrados del Error ($\text{SCE}$), inflando artificialmente el Cuadrado Medio del Error ($\text{CME}$) y reduciendo drásticamente la capacidad estadística del estadístico $F$ para detectar los efectos reales de los tratamientos.

El \textbf{Diseño en Bloques Completos al Azar (DBCA)} resuelve este problema agrupando las unidades experimentales en $b$ \textbf{bloques homogéneos} y aplicando los $k$ tratamientos aleatoriamente dentro de cada bloque, de modo que cada tratamiento aparece exactamente una vez por bloque ($N = k \times b$).

\subsection{Modelo aditivo y partición bidimensional de la varianza}

El modelo estadístico paramétrico para un diseño en bloques completos al azar se escribe como:
\begin{align}
	Y_{ij} = \mu + \tau_i + \beta_j + \epsilon_{ij}, \quad \text{para } i = 1, \ldots, k \text{ y } j = 1, \ldots, b,
	\label{eq:modelo_dbca}
\end{align}
donde:
\begin{itemize}
	\item $\mu$ es el promedio global de la población.
	\item $\tau_i$ es el efecto principal del \textbf{tratamiento $i$}, sujeto a la restricción $\sum_{i=1}^k \tau_i = 0$.
	\item $\beta_j$ es el efecto principal del \textbf{bloque $j$}, sujeto a la restricción $\sum_{j=1}^b \beta_j = 0$.
	\item $\epsilon_{ij} \sim N(0, \sigma^2)$ son los errores experimentales aleatorios i.i.d., asumiendo además la \textbf{ausencia de interacción} (aditividad pura) entre tratamientos y bloques.
\end{itemize}

\begin{teorema}[Teorema de partición de la suma de cuadrados en DBCA]
	En un Diseño en Bloques Completos al Azar con $k$ tratamientos y $b$ bloques, la Suma de Cuadrados Total se descompone en tres fuentes ortogonales de variación:
	\begin{align}
		\text{SCT} = \text{SCTR} + \text{SCB} + \text{SCE},
	\end{align}
	donde:
	\begin{itemize}
		\item $\text{SCTR} = b \sum_{i=1}^k (\bar{Y}_{i.} - \bar{Y}_{..})^2$ es la Suma de Cuadrados de los Tratamientos ($\nu_{\text{TR}} = k - 1$).
		\item $\text{SCB} = k \sum_{j=1}^b (\bar{Y}_{.j} - \bar{Y}_{..})^2$ es la \textbf{Suma de Cuadrados de los Bloques} ($\nu_{\text{B}} = b - 1$), donde $\bar{Y}_{.j} = \frac{1}{k}\sum_{i=1}^k Y_{ij}$ es el promedio observado de todo el bloque $j$.
		\item $\text{SCE} = \sum_{i=1}^k \sum_{j=1}^b (Y_{ij} - \bar{Y}_{i.} - \bar{Y}_{.j} + \bar{Y}_{..})^2$ es la Suma de Cuadrados del Error aleatorio residual, calculada en la práctica por diferencia algebraica directa:
		\begin{align}
			\text{SCE} = \text{SCT} - \text{SCTR} - \text{SCB},
		\end{align}
		con sus respectivos grados de libertad $\nu_E = (k - 1)(b - 1)$.
	\end{itemize}
\end{teorema}

El DBCA permite evaluar dos hipótesis nulas independientes: la prueba primordial sobre la igualdad de tratamientos ($H_{0,\text{TR}}: \tau_i = 0$) y la prueba secundaria para verificar si el factor de bloqueo redujo efectivamente el ruido del experimento ($H_{0,\text{B}}: \beta_j = 0$). Los estadísticos y su síntesis computacional se presentan en la Tabla \ref{tab:anova_dbca}.

\begin{table}[h!]
	\centering
	\begin{tabular}{lcccc}
		\hline
		\textbf{Fuente de Variación} & \textbf{SC} & \textbf{g.l. ($\nu$)} & \textbf{Cuadrado Medio (CM)} & \textbf{Estadístico $F$} \\ \hline
		Tratamientos & $\text{SCTR} = b \sum (\bar{Y}_{i.} - \bar{Y}_{..})^2$ & $k - 1$ & $\text{CMTR} = \frac{\text{SCTR}}{k - 1}$ & $F_{\text{TR}} = \frac{\text{CMTR}}{\text{CME}}$ \\ [1.5ex]
		Bloques & $\text{SCB} = k \sum (\bar{Y}_{.j} - \bar{Y}_{..})^2$ & $b - 1$ & $\text{CMB} = \frac{\text{SCB}}{b - 1}$ & $F_{\text{B}} = \frac{\text{CMB}}{\text{CME}}$ \\ [1.5ex]
		Error residual & $\text{SCE} = \text{SCT} - \text{SCTR} - \text{SCB}$ & $(k-1)(b-1)$ & $\text{CME} = \frac{\text{SCE}}{(k-1)(b-1)}$ & \\ [1.5ex] \hline
		\textbf{Total} & $\text{SCT} = \sum \sum (Y_{ij} - \bar{Y}_{..})^2$ & $kb - 1$ & & \\ \hline
	\end{tabular}
	\caption{Tabla de Análisis de Varianza para un Diseño en Bloques Completos al Azar (DBCA).}
	\label{tab:anova_dbca}
\end{table}

\subsection{Cuadrados latinos}

El DBCA controla la variabilidad proveniente de \textbf{una sola} fuente de bloqueo conocida. Cuando existen \textbf{dos} fuentes de variabilidad extrañas conocidas y cruzadas entre sí (por ejemplo, el turno de trabajo y la máquina utilizada en una línea de producción), el \textbf{diseño de cuadrado latino} permite controlar ambas simultáneamente sin que el número de corridas experimentales crezca multiplicativamente.

\begin{definicion}[Cuadrado latino]
	Un \emph{cuadrado latino} de orden $n$ es un arreglo de $n \times n$ celdas en el que se distribuyen $n$ tratamientos (denotados tradicionalmente con letras latinas $A, B, C, \ldots$) de manera que cada tratamiento aparece \textbf{exactamente una vez en cada fila y exactamente una vez en cada columna}. Las filas y columnas representan los dos factores de bloqueo; el número de corridas experimentales es $n^2$ en lugar de las $n^3$ que requeriría un diseño factorial completo con tres factores de $n$ niveles cada uno.
\end{definicion}

\begin{center}
	\begin{tabular}{c|cccc}
		 & Columna 1 & Columna 2 & Columna 3 & Columna 4 \\ \hline
		Fila 1 & $A$ & $B$ & $C$ & $D$ \\
		Fila 2 & $B$ & $C$ & $D$ & $A$ \\
		Fila 3 & $C$ & $D$ & $A$ & $B$ \\
		Fila 4 & $D$ & $A$ & $B$ & $C$ \\
	\end{tabular}
\end{center}

El modelo estadístico paramétrico para un cuadrado latino de orden $n$ es
\begin{align}
	Y_{ijk} = \mu + \rho_i + \gamma_j + \tau_k + \epsilon_{ijk},
	\label{eq:modelo_cuadrado_latino}
\end{align}
donde $\rho_i$ es el efecto de la fila $i$, $\gamma_j$ el efecto de la columna $j$, y $\tau_k$ el efecto del tratamiento $k$ (con $i,j,k = 1,\ldots,n$, sujetos a las restricciones usuales $\sum \rho_i = \sum \gamma_j = \sum \tau_k = 0$); nótese que, a diferencia del DBCA, solo se observa \emph{una} de las $n^2$ combinaciones fila-columna-tratamiento posibles por cada par $(i,j)$, ya que el propio arreglo del cuadrado latino determina qué tratamiento $k$ corresponde a cada celda.

\begin{teorema}[Partición de la suma de cuadrados en un cuadrado latino]
	La Suma de Cuadrados Total se descompone en cuatro fuentes ortogonales de variación:
	\begin{align}
		\text{SCT} = \text{SC}_{\text{filas}} + \text{SC}_{\text{columnas}} + \text{SCTR} + \text{SCE},
	\end{align}
	con $n^2 - 1$ grados de libertad totales, repartidos como $\nu_{\text{filas}} = \nu_{\text{columnas}} = \nu_{\text{TR}} = n-1$ y $\nu_E = (n-1)(n-2)$ para el error residual.
\end{teorema}

\begin{center}
	\begin{tabular}{lcccc}
		\hline
		\textbf{Fuente de Variación} & \textbf{SC} & \textbf{g.l. ($\nu$)} & \textbf{Cuadrado Medio (CM)} & \textbf{Estadístico $F$} \\ \hline
		Filas & $\text{SC}_{\text{filas}}$ & $n-1$ & $\text{SC}_{\text{filas}}/(n-1)$ & --- \\ [1ex]
		Columnas & $\text{SC}_{\text{columnas}}$ & $n-1$ & $\text{SC}_{\text{columnas}}/(n-1)$ & --- \\ [1ex]
		Tratamientos & $\text{SCTR}$ & $n-1$ & $\text{CMTR} = \frac{\text{SCTR}}{n-1}$ & $F_{\text{TR}} = \frac{\text{CMTR}}{\text{CME}}$ \\ [1ex]
		Error residual & $\text{SCE}$ & $(n-1)(n-2)$ & $\text{CME} = \frac{\text{SCE}}{(n-1)(n-2)}$ & \\ [1ex] \hline
		\textbf{Total} & $\text{SCT}$ & $n^2 - 1$ & & \\ \hline
	\end{tabular}
\end{center}

\begin{observacion}
	Al igual que en el DBCA, la hipótesis de interés primordial es $H_{0,\text{TR}}: \tau_k = 0$ para todo $k$ (igualdad de tratamientos), contrastada mediante $F_{\text{TR}}$. Las filas y columnas son bloques de control, no variables de interés experimental directo, por lo que usualmente no se reportan sus estadísticos $F$.
\end{observacion}

\subsection{Cuadrados grecolatinos}

Cuando existe una \textbf{tercera} fuente de variabilidad extraña conocida (además de las dos ya controladas por filas y columnas), el \textbf{cuadrado grecolatino} extiende la idea del cuadrado latino superponiendo un segundo arreglo de letras (tradicionalmente griegas: $\alpha, \beta, \gamma, \ldots$) sobre el mismo cuadrado $n \times n$.

\begin{definicion}[Cuadrado grecolatino]
	Un \emph{cuadrado grecolatino} de orden $n$ consiste en dos cuadrados latinos \textbf{ortogonales} superpuestos: uno con letras latinas ($A, B, C, \ldots$, el factor de tratamiento principal) y otro con letras griegas ($\alpha, \beta, \gamma, \ldots$, un cuarto factor controlado), de modo que cada combinación (letra latina, letra griega) aparece exactamente una vez en todo el arreglo. Se dice que los dos cuadrados son \emph{ortogonales} porque cada par (letra latina, letra griega) coincide en exactamente una celda.
\end{definicion}

\begin{observacion}[Condición de existencia]
	No existen cuadrados grecolatinos ortogonales de orden $n=6$ (resultado célebre conjeturado por Euler en 1782 y probado en 1900); para todo otro $n \geq 3$ sí existen. En la práctica, los cuadrados grecolatinos se usan casi exclusivamente para $n \geq 3$, $n \neq 6$.
\end{observacion}

El modelo estadístico paramétrico añade un cuarto término al modelo del cuadrado latino:
\begin{align}
	Y_{ijkl} = \mu + \rho_i + \gamma_j + \tau_k + \delta_l + \epsilon_{ijkl},
	\label{eq:modelo_cuadrado_grecolatino}
\end{align}
donde $\delta_l$ es el efecto del factor griego $l$ (con $i,j,k,l = 1,\ldots,n$). La Suma de Cuadrados Total se descompone ahora en cinco fuentes ortogonales:
\begin{align}
	\text{SCT} = \text{SC}_{\text{filas}} + \text{SC}_{\text{columnas}} + \text{SCTR} + \text{SC}_{\text{griego}} + \text{SCE},
\end{align}
con $\nu_{\text{filas}} = \nu_{\text{columnas}} = \nu_{\text{TR}} = \nu_{\text{griego}} = n-1$ y $\nu_E = (n-1)(n-3)$ grados de libertad para el error residual, sobre un total de $n^2 - 1$.

\begin{observacion}
	Cada fuente de bloqueo adicional (fila, columna, factor griego) que efectivamente explica variabilidad real reduce el Cuadrado Medio del Error y, por lo tanto, aumenta la potencia de la prueba $F$ sobre los tratamientos —el mismo principio de bloqueo del DBCA, aplicado ahora a tres fuentes de ruido simultáneas en lugar de una sola.
\end{observacion}

\section{Introducción a diseños factoriales}

Todos los diseños estudiados hasta ahora en este capítulo estudian el efecto de un \textbf{único} factor de tratamiento (posiblemente controlando fuentes de ruido adicionales mediante bloqueo). Sin embargo, muchas preguntas experimentales reales involucran \textbf{dos o más factores de tratamiento simultáneos} cuyo efecto conjunto —no solo individual— interesa evaluar: por ejemplo, el efecto combinado de la temperatura y la presión sobre el rendimiento de una reacción química, o el efecto conjunto del algoritmo y el tamaño del lote de entrenamiento sobre la exactitud de un modelo de aprendizaje automático.

\begin{definicion}[Diseño factorial completo]
	Un \emph{diseño factorial completo} con dos factores $A$ (con $a$ niveles) y $B$ (con $b$ niveles) consiste en observar \textbf{todas} las $a \times b$ combinaciones posibles de niveles de ambos factores, con $n \geq 1$ réplicas independientes por combinación (para un total de $N = abn$ observaciones). La generalización a más de dos factores sigue el mismo principio: se cruzan todos los niveles de todos los factores.
\end{definicion}

\begin{observacion}
	Frente a la alternativa de variar un factor a la vez manteniendo los demás fijos (\emph{one-factor-at-a-time}), el diseño factorial completo es estrictamente más eficiente: con el mismo número de corridas experimentales, permite estimar no solo los efectos principales de cada factor, sino también sus \textbf{interacciones}, algo que el enfoque secuencial no puede detectar.
\end{observacion}

\subsection{Efectos principales e interacción}

\begin{definicion}[Efecto principal e interacción]
	El \emph{efecto principal} de un factor es el cambio promedio en la variable respuesta al pasar de un nivel del factor a otro, promediando sobre todos los niveles del(los) otro(s) factor(es). Existe \emph{interacción} entre dos factores $A$ y $B$ cuando el efecto de $A$ sobre la variable respuesta \textbf{depende del nivel} de $B$ en el que se evalúe (y viceversa); es decir, las curvas de respuesta de $A$ para distintos niveles de $B$ no son paralelas.
\end{definicion}

El modelo estadístico paramétrico para un diseño factorial completo de dos factores es
\begin{align}
	Y_{ijk} = \mu + \alpha_i + \beta_j + (\alpha\beta)_{ij} + \epsilon_{ijk}, \quad i=1,\ldots,a,\ j=1,\ldots,b,\ k=1,\ldots,n,
	\label{eq:modelo_factorial}
\end{align}
donde $\alpha_i$ y $\beta_j$ son los efectos principales de $A$ y $B$ respectivamente, y $(\alpha\beta)_{ij}$ es el \textbf{efecto de interacción} del par de niveles $(i,j)$, sujeto a las restricciones usuales $\sum_i \alpha_i = \sum_j \beta_j = \sum_i (\alpha\beta)_{ij} = \sum_j (\alpha\beta)_{ij} = 0$.

\begin{teorema}[Partición de la suma de cuadrados en un diseño factorial de dos factores]
	La Suma de Cuadrados Total se descompone en cuatro fuentes ortogonales de variación:
	\begin{align}
		\text{SCT} = \text{SC}_A + \text{SC}_B + \text{SC}_{AB} + \text{SCE},
	\end{align}
	con grados de libertad $\nu_A = a-1$, $\nu_B = b-1$, $\nu_{AB} = (a-1)(b-1)$ para la interacción, y $\nu_E = ab(n-1)$ para el error residual, sobre un total de $\nu_T = abn - 1$.
\end{teorema}

\begin{center}
	\begin{tabular}{lcccc}
		\hline
		\textbf{Fuente de Variación} & \textbf{SC} & \textbf{g.l. ($\nu$)} & \textbf{Cuadrado Medio (CM)} & \textbf{Estadístico $F$} \\ \hline
		Factor $A$ & $\text{SC}_A$ & $a-1$ & $\text{CM}_A = \frac{\text{SC}_A}{a-1}$ & $F_A = \frac{\text{CM}_A}{\text{CME}}$ \\ [1ex]
		Factor $B$ & $\text{SC}_B$ & $b-1$ & $\text{CM}_B = \frac{\text{SC}_B}{b-1}$ & $F_B = \frac{\text{CM}_B}{\text{CME}}$ \\ [1ex]
		Interacción $AB$ & $\text{SC}_{AB}$ & $(a-1)(b-1)$ & $\text{CM}_{AB} = \frac{\text{SC}_{AB}}{(a-1)(b-1)}$ & $F_{AB} = \frac{\text{CM}_{AB}}{\text{CME}}$ \\ [1ex]
		Error residual & $\text{SCE}$ & $ab(n-1)$ & $\text{CME} = \frac{\text{SCE}}{ab(n-1)}$ & \\ [1ex] \hline
		\textbf{Total} & $\text{SCT}$ & $abn-1$ & & \\ \hline
	\end{tabular}
\end{center}

\begin{observacion}[Orden de interpretación]
	Por convención, un análisis factorial siempre se interpreta \textbf{de adentro hacia afuera}: se examina primero el estadístico $F_{AB}$ de la interacción. Si resulta significativo, los efectos principales $\alpha_i$ y $\beta_j$ ya no tienen una interpretación aislada simple (el efecto de $A$ cambia según el nivel de $B$), y el análisis debe reportarse en términos de las combinaciones específicas de niveles. Solo cuando la interacción \textbf{no} es significativa tiene sentido interpretar directamente los efectos principales por separado.
\end{observacion}

\begin{ejemplo}
	\label{exmp:8.6.1}
	Un ingeniero de datos evalúa el tiempo de entrenamiento (en minutos) de un modelo de aprendizaje automático bajo dos algoritmos de optimización ($A_1$, $A_2$) y dos tamaños de lote ($B_1 = 32$, $B_2 = 256$), con $n=3$ réplicas por combinación. Las medias muestrales observadas por celda son:
	\begin{center}
		\begin{tabular}{lcc}
			\hline
			 & $B_1$ (lote 32) & $B_2$ (lote 256) \\ \hline
			$A_1$ & $42$ & $40$ \\
			$A_2$ & $38$ & $20$ \\ \hline
		\end{tabular}
	\end{center}
	¿Sugieren estos promedios la presencia de interacción entre el algoritmo y el tamaño de lote?
\end{ejemplo}

\begin{solucion}
	El efecto de cambiar de $B_1$ a $B_2$ es muy distinto según el algoritmo: bajo $A_1$, el tiempo apenas cambia ($42 \to 40$, una reducción de $2$ minutos), mientras que bajo $A_2$ el tiempo se reduce drásticamente ($38 \to 20$, una reducción de $18$ minutos). Como el efecto del factor $B$ (tamaño de lote) depende marcadamente del nivel del factor $A$ (algoritmo), los datos sugieren una \textbf{interacción sustancial} entre ambos factores. Esto significa que no es correcto afirmar de forma general ``el lote grande reduce el tiempo de entrenamiento''; la afirmación correcta y útil es que el lote grande reduce el tiempo \emph{solo bajo el algoritmo $A_2$}. Confirmar formalmente esta sospecha requeriría calcular $F_{AB}$ y compararlo contra el cuantil crítico $F_{\alpha,\,(a-1)(b-1),\,ab(n-1)}$.
\end{solucion}
